\exercisesection{\S~10}

\begin{enumerate}
\def\labelenumi{\arabic{enumi})}
\item
  Let \(\mathfrak{s} = \mathfrak{sl}(2, k)\) and \(\mathfrak{g} = \mathfrak{sl}(3, k)\) . Identify \(\mathfrak{s}\) with a subalgebra of \(\mathfrak{g}\) by means of an irreducible representation of \(\mathfrak{s}\) of degree 3. Show that every subspace of \(\mathfrak{g}\) containing \(\mathfrak{s}\) and stable under \(\operatorname{ad}_{\mathfrak{g}}\mathfrak{s}\) is equal to either \(\mathfrak{s}\) or \(\mathfrak{g}\) ; deduce that \(\mathfrak{s}\) is maximal among the subalgebras of \(\mathfrak{g}\) distinct from \(\mathfrak{g}\) .
\item
  Let \(m = \frac{1}{2}(\dim(\mathfrak{g}) + \operatorname{rk}(\mathfrak{g}))\) . Every solvable subalgebra of g is of dimension \(\leq m\) ; if it is of dimension m, it is a Borel subalgebra. (Reduce to the algebraically closed case, and use Th. 2.)
\item
  Assume that \(k\) is \(\mathbf{R}, \mathbf{C}\) , or a non-discrete complete ultrametric field. Give the grassmannian \(\mathbf{G}(\mathfrak{g})\) of vector subspaces of \(\mathfrak{g}\) its natural structure of analytic manifold over \(k\) (Differentiable and Analytic Manifolds, Results, 5.2.6). Consider the subsets of \(\mathbf{G}(\mathfrak{g})\) formed by:
\end{enumerate}

\begin{enumerate}
\def\labelenumi{(\roman{enumi})}
\item
  the subalgebras
\item
  the solvable subalgebras
\item
  the nilpotent subalgebras
\item
  the subalgebras consisting of nilpotent elements
\item
  the Borel subalgebras.
\end{enumerate}

Show that these subsets are closed (for (v), use Exerc. 2). Deduce that these subsets are compact when \(k\) is locally compact.

Show by examples that the subsets of \(\mathbf{G}(\mathfrak{g})\) formed by:

\begin{enumerate}
\def\labelenumi{(\roman{enumi})}
\setcounter{enumi}{5}
\item
  the Cartan subalgebras
\item
  the subalgebras reductive in \(\mathfrak{g}\)
\item
  the semi-simple subalgebras
\item
  the decomposable subalgebras
\end{enumerate}

are not necessarily closed, even when \(k = \mathbf{C}\) .

\begin{enumerate}
\def\labelenumi{\arabic{enumi})}
\setcounter{enumi}{3}
\tightlist
\item
  Assume that \(k = \mathbf{C}\) . Let \(G = \operatorname{Int}(\mathfrak{g}) = \operatorname{Aut}_0(\mathfrak{g})\) , and let \(B\) be an integral subgroup of \(G\) whose Lie algebra \(\mathfrak{b}\) is a Borel subalgebra of \(\mathfrak{g}\) . Show that \(B\) is the normalizer of \(\mathfrak{b}\) in \(G\) (use Exerc. 4 of §3 and Exerc. 11 of §5). Deduce by means of Exerc. 3 that \(G / B\) is compact.
\end{enumerate}

¶5) Assume that g is splittable. If h is a splitting Cartan subalgebra of g, denote by \(\mathrm{E}(\mathfrak{h})\) the subgroup of \(\mathrm{Aut}_{e}(\mathfrak{g})\) generated by the \(e^{\operatorname{ad} x}, x \in \mathfrak{g}^{\alpha}(\mathfrak{h}), \alpha \in \mathrm{R}(\mathfrak{g}, \mathfrak{h})\) , cf.~Chap. VII, §3, no. 2.

\begin{enumerate}
\def\labelenumi{\alph{enumi})}
\item
  Let \(\mathfrak{b}\) be a Borel subalgebra of \(\mathfrak{g}\) containing \(\mathfrak{h}\) , and let \(\mathfrak{h}_1\) be a Cartan subalgebra of \(\mathfrak{b}\) . Show that \(\mathfrak{h}_1\) is conjugate to \(\mathfrak{h}\) by an element of \(\mathrm{E}(\mathfrak{h})\) . Deduce that \(\mathrm{E}(\mathfrak{h}) = \mathrm{E}(\mathfrak{h}_1)\) .
\item
  Let \(\mathfrak{h}'\) be a splitting Cartan subalgebra of \(\mathfrak{g}\) . Show that \(\mathrm{E}(\mathfrak{g}) = \mathrm{E}(\mathfrak{h}')\) . (If \(\mathfrak{b}'\) is a Borel subalgebra containing \(\mathfrak{h}'\) , choose a Cartan subalgebra \(\mathfrak{h}_1\) of \(\mathfrak{b} \cap \mathfrak{b}'\) and apply \(a\) ) to show that \(\mathrm{E}(\mathfrak{h}) = \mathrm{E}(\mathfrak{h}_1) = \mathrm{E}(\mathfrak{h}')\) .)
\item
  Let \(x\) be a nilpotent element of \(\mathfrak{g}\) . Show that \(e^{\mathrm{ad}x} \in \mathrm{E}(\mathfrak{h})\) . (Use b) and Cor. 2 of Th. 1 to reduce to the case in which \(x \in [\mathfrak{b}, \mathfrak{b}]\) .) Deduce that \(\mathrm{E}(\mathfrak{h}) = \mathrm{Aut}_e(\mathfrak{g})\) .
\end{enumerate}

\begin{enumerate}
\def\labelenumi{\arabic{enumi})}
\setcounter{enumi}{5}
\tightlist
\item
  \begin{enumerate}
  \def\labelenumii{\alph{enumii})}
  \tightlist
  \item
    Show the equivalence of the properties:
  \end{enumerate}
\end{enumerate}

\begin{enumerate}
\def\labelenumi{(\roman{enumi})}
\item
  g has no nilpotent element \(\neq 0\) .
\item
  g has no parabolic subalgebra ≠ g.
\end{enumerate}

(Use Cor. 2 of Th. 1.)

Such an algebra is called anisotropic.

\begin{enumerate}
\def\labelenumi{\alph{enumi})}
\setcounter{enumi}{1}
\tightlist
\item
  Let \(\mathfrak{p}\) be a minimal parabolic subalgebra of \(\mathfrak{g}\) , \(\mathfrak{r}\) the radical of \(\mathfrak{p}\) , and \(\mathfrak{s} = \mathfrak{p} / \mathfrak{r}\) . Show that \(\mathfrak{s}\) is anisotropic. (Remark that, if \(\mathfrak{q}\) is a parabolic subalgebra of \(\mathfrak{s}\) , the inverse image of \(\mathfrak{q}\) in \(\mathfrak{p}\) is a parabolic subalgebra of \(\mathfrak{g}\) , cf.~§3, Exerc. 5 a).)
\end{enumerate}

¶ 7) a) Show that the following properties of k are equivalent:

\begin{enumerate}
\def\labelenumi{(\roman{enumi})}
\item
  Every anisotropic semi-simple \(k\) -Lie algebra reduces to 0.
\item
  Every semi-simple k-Lie algebra has a Borel subalgebra.
\end{enumerate}

(Use Exerc. 6 to prove that (i) \(\Longrightarrow\) (ii).)

\begin{enumerate}
\def\labelenumi{\alph{enumi})}
\setcounter{enumi}{1}
\tightlist
\item
  Show that (i) and (ii) imply \(^{19}\) :
\end{enumerate}

\begin{enumerate}
\def\labelenumi{(\roman{enumi})}
\setcounter{enumi}{2}
\tightlist
\item
  Every finite dimensional \(k\) -algebra that is a field is commutative. (Or again: the Brauer group of every algebraic extension of \(k\) reduces to 0.)
\end{enumerate}

(Use the Lie algebra of elements of trace zero in such an algebra.)

\begin{enumerate}
\def\labelenumi{\alph{enumi})}
\setcounter{enumi}{2}
\tightlist
\item
  Show that (i) and (ii) are implied by:
\end{enumerate}

\begin{enumerate}
\def\labelenumi{(\roman{enumi})}
\setcounter{enumi}{3}
\tightlist
\item
  For any finite family of homogeneous polynomials \(f_{\alpha} \in k[(X_i)_{i \in \mathrm{I}}]\) of degrees \(\geq 1\) such that \(\sum_{\alpha} \deg f_{\alpha} < \operatorname{Card}(\mathrm{I})\) , there exist elements \(x_i \in k\) , not all zero, such that \(f_{\alpha}((x_i)_{i \in \mathrm{I}}) = 0\) for all \(\alpha\) .
\end{enumerate}

(Use Prop. 5 of §8.)

